\clearpage\begingroup\SourceLeaf{118b}\HistoricalHead{113}{R}{GAUSS AN OLBERS UND ENCKE.}\HistoricalFoot{XI1.}{15}
\begin{center}[28.]\end{center}
\begin{center}\textsc{Gauss} an \textsc{Encke}.\quad Göttingen, 8. September 1839.\end{center}
\noindent .\ .\ .\ .\ .\ .\ .\ .\ .\ .\ .\ .\ .\ .\ .\ .\ .\ .\ .\ .\ .\par
\indent Das gütige Interesse, welches Sie an meiner allgemeinen Theorie des
Erdmagnetismus bezeigen, ist mir sehr erfreulich gewesen. Wenn diese Ar-
\clearpage\SourceLeaf{119}\HistoricalHead{114}{L}{PHYSIK. BRIEFWECHSEL.}\HistoricalFoot{}{}
\noindent beit ein Verdienst hat, so besteht es meiner Meinung nach darin, dass man
früher nicht wusste, wie man überhaupt die Sache anzufassen habe, und dass
jetzt der Weg dazu offen ist. Ich habe insofern eine numerische Anwendung
ungern gemacht, als ich die Ueberzeugung hatte und habe, dass die groben
Data jetzt eine so grosse mechanische Arbeit, wie ich habe aufwenden müssen,
kaum verdienen, und dass bei guten Datis eine \textls[120]{nicht längere} Arbeit ganz
anders befriedigende Resultate hätte liefern können und müssen. Indessen
war ich eben so sehr überzeugt, dass eine \textls[120]{blosse} Auseinandersetzung der
Theorie gar keinen Eingang gefunden haben würde.

In Beziehung auf Ein Moment scheinen Sie übrigens sich nicht ganz in
den eigentlichen Standpunkt gesetzt zu haben, nämlich dass Sie pag. 47\textsuperscript{[*)]} einen
\textls[120]{Anstoss} gefunden haben. Sie scheinen dieses nur wie einen kleinen Nebenpunkt
zu betrachten, wo Ihnen nicht gleich gelungen sei die Lücke zu suppliren.
Es ist aber gar nicht meine Meinung gewesen, dass die Sache aus
\textls[120]{diesem} Gesichtspunkt betrachtet werden soll. Vielmehr ist diess nur Ein
Hauptsatz aus einer grösseren theoretischen Untersuchung, auf die ich selbst
nach meinem eigenen Maassstabe, einen viel grösseren Werth lege, als auf
die ganze Abhandlung, die jetzt gedruckt ist, insofern in \textls[120]{jener} Arbeit 10mahl
mehr \textls[120]{geistige von einander unabhängige} Schritte zu machen sind, als
in dieser, wo eigentlich nach der einmal gefassten Grund-Idee alles sich von
selbst macht.\quad .\ .\ .\ .\ .\ .\ .\ .\ .\ .\ .\ .\ .\ .\ .\ .\ .\ .\ .\ .\ .
\begingroup\renewcommand{\thefootnote}{}\footnotetext{[*) Es handelt sich genauer um Seite 46 und 47 der Resultate aus den Beobachtungen des Magnetischen Vereins im Jahre 1838, wo im Artikel 32 der \textit{Allgemeinen Theorie des Erdmagnetismus} (siehe Werke V, S. 162) das Theorem des art. 2. der \textit{Intensitas vis magneticae terrestris} erwähnt wird, vergl. oben S. 103 die zweite Fussnote.]}\endgroup
\clearpage\endgroup
